% TOP_PROBLEM: {"id": 1513, "title": "Four Exponentials Conjecture", "category_id": 1, "category": {"id": 1, "name": "number_theory", "display_name": "Number Theory", "description": "Properties of integers, prime numbers, Diophantine equations.", "slug": "number-theory", "order_index": 1, "created_at": "2026-07-31T15:26:25.670Z"}, "status": "open"}
% FIRST_POSTED: null
% ENTRY_KIND: statement_only
% Imported from: batch08.tex; UnsolvedMath ID 1513
\documentclass[11pt]{article}
\usepackage[margin=25mm]{geometry}
\usepackage{fontspec}
\setmainfont{Latin Modern Roman}
\usepackage{amsmath,amssymb,mathtools}
\usepackage{unicode-math}
\setmathfont{Latin Modern Math}
\usepackage{xurl,hyperref}
\hypersetup{colorlinks=true,urlcolor=blue}
\setcounter{secnumdepth}{0}
\setlength{\parindent}{0pt}
\setlength{\parskip}{5pt}
\setlength{\emergencystretch}{3em}
\newcommand{\sref}[2]{\hyperlink{source:#1:#2}{[S#2]}}
\newcommand{\cataloguescope}[1]{\paragraph{Recorded target.}#1}
\begin{document}
% UnsolvedMath ID 1513; source catalogue code NUM-019
\setcounter{section}{1512}
\section{Four Exponentials Conjecture}\label{1513}
{\small\sffamily UnsolvedMath ID 1513 · Number Theory}
\subsection{Definitions and mathematical statement}

\paragraph{Shared notation.}$\mathbb N=\{1,2,\ldots\}$, and $\mathbb Z,\mathbb Q,\mathbb R,\mathbb C$ denote the integers, rationals, reals and complex numbers. Any other conventions are specified locally.


Two complex numbers \(u_1,u_2\) are linearly independent over
\(\mathbb Q\) if \(a_1u_1+a_2u_2=0\), with
\(a_1,a_2\in\mathbb Q\), implies \(a_1=a_2=0\).
A complex number is algebraic if it is a root of a nonzero polynomial
with rational coefficients, and transcendental otherwise.
For \(z\in\mathbb C\), \(e^z=\sum_{m\ge0}z^m/m!\) is the
usual complex exponential.
\paragraph{Statement recorded in the supplied source.}
For every pair \(x_1,x_2\in\mathbb C\) linearly independent over
\(\mathbb Q\), and every pair \(y_1,y_2\in\mathbb C\) linearly
independent over \(\mathbb Q\), is at least one of
\[
             e^{x_1y_1},\ e^{x_1y_2},\ e^{x_2y_1},\ e^{x_2y_2}
\]
transcendental? Independence is required separately within each pair. \sref{1513}{2}
\subsection{Short English statement}
Must a two-by-two array of exponentials contain a transcendental value when its two row parameters and its two column parameters are each rationally independent?
\subsection{Sources}
\begin{description}
\item[\hypertarget{source:1513:1}{S1}] ULAM.AI and the UnsolvedMath Contributors, \emph{UnsolvedMath: A Curated Collection of Open Mathematics Problems}, record 1513 (NUM-019). CC BY 4.0. \url{https://huggingface.co/datasets/ulamai/UnsolvedMath}.
\item[\hypertarget{source:1513:2}{S2}] Michel Waldschmidt, \emph{Open Diophantine Problems}, arXiv:math/0312440v2 (2004), Section 3.1, Conjecture 3.7 (Four Exponentials Conjecture). \url{https://arxiv.org/abs/math/0312440}.
\end{description}
\label{end:1513}
\end{document}
